%% 06_limitations.tex
% Numbers from: analysis/results.md, analysis/feature_extraction.md

\section{Limitations}
\label{sec:limitations}

\paragraph{Synthetic data only.}
All 200 patients are generated by Synthea with a custom specialty-regimen overlay.
Synthea's medication module is known to deviate from real prescribing
distributions~\citep{hodges2023medication} and models only guideline-concordant
pathways without real-world deviations~\citep{meeker2022synthea}.  Our overlay
addresses the specialty-medication gap but cannot introduce the messy, noisy,
protocol-discordant adherence patterns present in real Patient Support Program data.
Absolute accuracy numbers should not be extrapolated to production clinical workflows.
External validity to real de-identified PSP data is the explicit follow-up step.

\paragraph{ANSWER\_MAX\_TOKENS = 512 is a partial confounder.}
The 512-token output budget contributes to 44--58\% of sampled failures via
reasoning truncation (Table~\ref{tab:errors}) and disproportionately penalises
structured systems whose FHIR-JSON contexts require longer chain-of-thought
arithmetic.  The 1024-token sensitivity sweep (\S\ref{sec:results:budget_sweep})
quantifies this contribution on 495 paired questions: doubling the budget gives
$+0$~pp to System~A, $+3.8$~pp to System~B, and $+5.1$~pp to System~C, narrowing
but not eliminating A's lead.  The A $>$ B $>$ C ordering survives at 1024 tokens.
A full sweep across additional budget points (256, 768, 2048) is left for the
full-venue extension.

\paragraph{Single answer LLM and single embedding model.}
Results are conditioned on \path{qwen/qwen3-32b} (Groq) as the answer LLM and
\path{BAAI/bge-large-en-v1.5} as the embedding model.  Robustness across other
LLMs or domain-adapted embeddings (MedCPT~\citep{jin2023medcpt}) is flagged as future
work.  The templated-narrative ablation (deterministic generator, no LLM) was generated
and passed the same fidelity audit at 100\% entity recall across all 200 patients and
all entity classes.  A downstream QA evaluation against the templated narratives was
planned to test robustness of the A $>$ B $>$ C ordering to narrative-LLM choice, but
was not completed in time for this submission.  Whether structured systems would beat
or tie templated narratives remains an open question, deferred to the full-venue version.

\paragraph{Per-cell N at tier $\times$ family intersections.}
With 200 patients across 3 tiers and 5 families, per-cell question counts range from
504 (Tier~1 $\times$ \texttt{cross\_resource}) to 1{,}554 (Tier~3 $\times$
\texttt{regimen\_compliance}).  Wilson CIs at those cell sizes are approximately
$\pm$3~pp; comparisons within cells should be treated as exploratory.

\paragraph{Single-rater error taxonomy.}
An ``Other / unclassified'' residual of 4--8\% indicates that the five-category taxonomy
(four primary plus residual) does not exhaust all failure modes.  Inter-rater reliability
is not reported because the design is single-rater by intent; the taxonomy itself is
deterministic and reproducible from the regex rules in
\texttt{analysis/run\_error\_taxonomy.py}.

\paragraph{Conservative retrieval recall heuristic.}
Recall@k for Systems~B and~C is computed using a heuristic that checks whether the
top-$k$ chunk IDs include the patient's primary MedicationRequest or specialty CarePlan.
This is a conservative proxy; a full per-question source-resource mapping (requiring
ground-truth functions in ``explain'' mode) is deferred to a follow-up release.

\paragraph{Feature-extraction arm: synthetic adherence label.}
The FS-Structured AUC near 1.0 reflects label/feature alignment in the synthetic design:
all positive cases are Tier-2 patients, and FS-Structured directly encodes tier as a
one-hot feature.  The near-perfect AUC should not be interpreted as evidence that
structured features will achieve similar performance on real, noisy adherence outcomes.
The feature-extraction arm is a design-informing benchmark, not a production predictor.

\paragraph{No zero-retrieval baseline.}
We do not report an arm in which the answer LLM is asked the same questions with no
retrieved context.  Such a baseline would partition the 40.6\% A accuracy into a
question-text-only component and a retrieval-attributable component.  The omission
limits direct attribution of the absolute accuracy numbers, but does not affect the
paired pairwise comparisons (A vs.\ B, A vs.\ C, B vs.\ C), all of which use the same
answer LLM and the same questions.

\paragraph{Class-imbalance handling for LightGBM.}
The synthetic adherence label is 14.5\% positive.  Logistic regression uses
\texttt{class\_weight=balanced}; LightGBM is trained with default class-weight settings
(no \texttt{is\_unbalance} or \texttt{scale\_pos\_weight} adjustment).  Because AUC-ROC
is threshold-independent, the reported AUCs are not biased by this choice; downstream
threshold-dependent metrics (precision, recall at a fixed cut-off) would be.

\paragraph{Dataset release contingent on IP clearance.}
Dataset release under CC-BY-4.0 is contingent on employer IP clearance and may be
delayed.  Code is released unconditionally under Apache-2.0.

\paragraph{Three tiers may not cover all real-world regimen complexity.}
The three complexity tiers are designed to span simple LAI, cyclic, and multi-drug
regimens, but do not cover all specialty-medication scenarios (infusion-only regimens,
combination-biologic prescribing, variable-interval dosing).  Adherence metrics
(PDC, MPR) are implemented to standard definitions~\citep{lociganic2015using} but edge
cases (overlapping refills, dose reductions, bridge therapies) are not represented in
the synthetic dataset.
